\section{Introduction}
\label{sec:intro}

A metadata-hiding group registry, in the sense we use the term, is a
public-state object that records the membership of a group as an
opaque commitment and admits state transitions only on production of a
zero-knowledge proof that some authorization predicate holds against
the prior state. Tornado Cash~\cite{pertsev-2019-tornado} fits this
shape with the predicate ``a member of the deposit set is unspending a
specific note''; Semaphore~\cite{semaphore-2024} fits it with ``a
member of an identity tree is signaling exactly once for a given
external nullifier''; Aztec~\cite{aztec-2024} fits it as the
application execution model rather than as a single contract; the
recent Stellar Ecosystem Proposal at~\cite{sep-onym} fits it with
group membership and per-epoch update authorization. The shape is
identifiable without committing to a curve, a SNARK system, or an
anchor chain.

The literature, however, treats the curve, SNARK, and chain as
bundled. A reader who wants to compare designs --- to ask whether a
given deployment should anchor on Stellar or on an Ethereum L2,
whether Groth16~\cite{groth-2016-onpairing} or Halo2 is more
appropriate for a given trust posture, whether the verifier should
ride a precompile or run in pure bytecode --- has no taxonomy to read
against. Existing surveys of zero-knowledge constructions index by
proof system or by hash family, not by the deployment constraints that
in practice determine which proof system is reachable. The
post-quantum migration question for this class of system is
particularly affected: it is not currently catalogued which migration
paths are technically blocked, which are economically infeasible, and
which are ready to ship.

This paper systematizes the design space.

\subsection*{What this paper is, and is not}

This is a Systematization of Knowledge paper. It is not a new
construction, not a new proof, and not an experience report. The
contribution is a taxonomy and a set of measurements taken against
that taxonomy, with one caveat about measurement scope:

\begin{itemize}
  \item For all configurations except one, we cite published
    benchmark numbers and date-stamp them. We do not re-measure
    third-party systems on uniform hardware; uniform-hardware
    re-measurement of every cell would be a separate paper, and we
    believe the available cited numbers are sufficient for the
    cross-configuration comparisons we draw.
  \item For the Stellar~+~BLS12-381~+~Groth16 row, we run a
    first-party benchmark harness (\S\ref{sec:analysis},
    appendix). One row of every comparison table is therefore
    authoritative under our own measurement protocol; the rest are
    cited.
\end{itemize}

We do not argue that any one configuration is the right configuration.
The recommendations in \S\ref{sec:recommendations} are conditional on
deployment profile --- a cost-optimized profile, a
post-quantum-prioritized profile, a censorship-resistance-prioritized
profile, and a jurisdictional-constraints-prioritized profile each
land on a different point in the design space.

\subsection*{Contributions}

\begin{enumerate}
  \item \textbf{A six-axis taxonomy} (\S\ref{sec:dimensions}): anchor
    chain, SNARK system, curve/hash pair, post-quantum stance, trusted
    setup model, and verification host. We argue that these axes are
    the right granularity for design discussion, decomposing the
    bundled (curve, SNARK, chain) framing that existing literature
    most often presents into its constituent decisions; cross-axis
    dependencies are made explicit rather than collapsed.
  \item \textbf{Per-dimension cost, trust, and post-quantum
    analysis} (\S\ref{sec:analysis}). Five tables: verification cost
    (in 2026-04 USD), proof size, prover wall-clock, trust
    assumptions, and post-quantum class.
  \item \textbf{Seven concrete configurations}
    (\S\ref{sec:configurations}) as labeled points in the taxonomy,
    including one ``hypothetical Stellar~+~Plonky3~+~FRI'' point that
    is currently blocked at the host-function layer. The blocked
    configuration is included precisely because it makes the
    host-function axis legible.
  \item \textbf{A migration-path catalogue} (\S\ref{sec:migration})
    classifying every starting-to-PQ-ready transformation as one of:
    ready today, blocked on host functions, blocked on tooling
    maturity, or blocked on cryptographic open problems.
  \item \textbf{Profile-driven deployment recommendations}
    (\S\ref{sec:recommendations}) that map each of four operational
    priorities to its dominant configuration in the taxonomy, with
    explicit non-recommendations for the cases where no configuration
    in the surveyed space is appropriate.
\end{enumerate}

\subsection*{Paper map}

\S\ref{sec:problem} states the problem class abstractly and fixes a
threat model. \S\ref{sec:dimensions} introduces the six design axes.
\S\ref{sec:analysis} populates the per-axis tables.
\S\ref{sec:configurations} catalogues seven concrete configurations.
\S\ref{sec:migration} surveys migration paths and what blocks them.
\S\ref{sec:recommendations} maps deployment profiles to
configurations. \S\ref{sec:open} enumerates open problems made visible
by the taxonomy. \S\ref{sec:related} positions the contribution
against prior work. \S\ref{sec:conclusion} concludes.
